% 04-moteur.tex — Le moteur : élan, lancer, alignement, lâcher.

\section{Le moteur : élan, lancer et lâcher}

Le seau part du repos, pendu sous la barre. Le moteur ne pousse jamais le seau
lui-même : il choisit l'angle d'avance \(\delta\) de la barre sur le seau, et
c'est la corde, tirée par l'extrémité de la barre, qui accélère le seau. Avec
\(r\) la distance pivot--seau, la corde fait alors l'angle \(\psi\) avec la
direction radiale du seau, donné par le triangle pivot--extrémité--seau :
\begin{equation}
  L_b \sin(\delta + \psi) = r \sin(\psi), \qquad \text{donc} \qquad \sin(\psi) \leq \frac{L_b}{r}.
\end{equation}
C'est pourquoi la barre est longue devant la corde : comme un bras qui fait
tourner une corde courte, elle peut incliner la traction.

\paragraph{Énergie apportée.} Par unité de masse, la corde doit fournir la
composante radiale \(r \dot\theta^2 + g \cos(\theta)\) qui maintient le seau
sur son cercle ; inclinée de \(\psi\), elle fournit en plus la composante
tangentielle \((r \dot\theta^2 + g \cos(\theta)) \tan(\psi)\). L'énergie
par unité de masse \(E = v^2/2 + g R (1 - \cos(\theta))\) varie donc, en
assimilant \(r\) à \(R\), selon
\begin{equation}
  \frac{\mathrm{d}E}{\mathrm{d}\theta} = \left(2E + g R \left(3 \cos(\theta) - 2\right)\right) \tan(\psi).
  \label{eq-energie}
\end{equation}
Le moteur inverse cette relation : il choisit \(\tan(\psi)\), puis \(\delta\),
pour que l'énergie suive une courbe voulue. La traction est limitée à
\(\tan(\psi) \leq 0{,}35\) pour que le seau, articulé à l'anse, suive.

\paragraph{Phases.} Le mouvement comporte quatre phases (\cref{fig-captures}).
\begin{itemize}
  \item \emph{Élan} : quelques allers-retours au rythme propre du pendule, en
        tirant dans le sens du mouvement, jusqu'à une amplitude d'environ
        \ang{70}.
  \item \emph{Lancer} : à partir d'un recul, le moteur fait suivre à
        l'énergie une courbe lisse qui atteint \(v_0^2/2\) à \ang{90}
        (équation~\eqref{eq-energie} plus une correction proportionnelle à
        l'écart).
  \item \emph{Alignement} : dès que l'énergie, augmentée de ce que
        l'alignement apportera encore, atteint \(v_0^2/2\), la barre revient
        dans l'axe de la corde et accompagne le seau jusqu'à ce que la chaîne
        ne plie plus (flexion inférieure à \ang{3}, au plus \SI{0,6}{s}).
  \item \emph{Lâcher} : la barre tourne ensuite librement, et le seau passe au
        point bas avec une vitesse proche de \(v_0\) (voir
        \cref{sec-validation}).
\end{itemize}
Un lâcher avec une chaîne pliée fait fouetter la barre : la corde se retend
brusquement et la tension filtrée atteint environ 2,3 fois sa valeur avant le
lâcher (\SI{230}{N} contre \SI{100}{N} dans un essai sur Terre, \(L = \SI{0,4}{m}\),
\(m = \SI{2}{kg}\)). C'est la raison de la phase d'alignement.
